% Volume VIII: Computational Neuroscience Mathematics
% Continuity note for Chapter 47.
% Previous dependencies: Chapter 13's random variables, joint distributions, random vectors, and covariance; Chapter 14's Poisson, Gaussian, categorical, and exponential-family models; Chapter 15's conditioning and Bayes' theorem; Chapter 18's entropy, cross-entropy, KL divergence, mutual information, data processing, and Fisher information.
% New durable concepts: tuning curves; preferred stimulus and tuning width; response variability; population response vectors; population vector decoding; Fisher information for population codes; Bayesian decoders; linear decoders; reconstruction loss; response entropy; noise entropy; stimulus-response mutual information; finite-sample bias and binning caveats.
% Forward hooks: Chapter 48 turns population responses into state-space geometry and latent trajectories; Chapter 49 studies recurrent circuits that generate population codes; Chapter 50 reuses Bayesian decoding as a model of perceptual inference; Chapter 51 compares encoding and decoding models on neural data.
% Notation commitments: Stimuli are S or s in a stimulus set \mathcal S; neural responses are R or r in a response set \mathcal R; neuron i has response R_i and tuning curve f_i(s)=\mathbb E[R_i\mid S=s]; population responses are R=(R_1,\ldots,R_N) in \mathbb R^N or \mathbb N^N; mutual information between stimulus and response is I(S;R).
\section{Chapter 47: Neural Coding, Decoding, and Information}
\label{ch:neural-coding-decoding-information}

The preceding chapters built probability, statistics, and information theory as general tools. Neural coding asks what those tools say about recordings from nervous systems. A stimulus, action, memory item, or latent decision variable is not observed directly inside a neuron. What we observe are spikes, calcium signals, voltages, local field potentials, or population activity. The mathematical problem is to relate those responses to the variables they may represent.

This chapter does not assume that neurons literally write messages in a private code. It uses ``code'' in the operational sense: a response distribution \(p(r\mid s)\) changes with the variable \(s\), and an observer or downstream circuit can use \(r\) to reduce uncertainty about \(s\). Encoding studies the forward map from variables to responses. Decoding studies the inverse map from responses to estimates. Information theory quantifies how much uncertainty is reduced, with finite-data warnings.

\paragraph{Chapter problem.}
How can tuning curves, population responses, decoders, and information measures make precise what neural activity says about stimuli, actions, or latent variables?

\paragraph{Section map.}
Section 47.1 defines tuning curves as conditional response averages and separates tuning from variability. Section 47.2 studies population codes as response vectors with geometry, covariance, and Fisher information. Section 47.3 derives Bayesian and linear decoders for reconstruction. Section 47.4 specializes entropy, KL divergence, cross-entropy, mutual information, and data processing to neural responses, including finite-data caveats.

\subsection{47.1 Tuning curves}

\paragraph{Motivating question.}
How does the expected response of one neuron change as a stimulus or task variable changes?

\paragraph{Beginner setup.}
A tuning curve is an input-output summary for one neuron. The input might be stimulus orientation, sound frequency, spatial location, reach direction, head direction, elapsed time, or a task variable. The output is usually a firing-rate estimate or spike-count average across repeated trials.

Suppose \(S\) is a stimulus and \(R_i\) is the response of neuron \(i\) in a fixed time window. The tuning curve of neuron \(i\) is
\[
f_i(s)=\mathbb E[R_i\mid S=s].
\]
If \(f_i(s)\) is largest near one value \(s^\star\), we call \(s^\star\) a preferred stimulus. If \(f_i(s)\) is broad, the neuron responds to many similar stimuli. If it is narrow, the neuron is selective.

The tuning curve is not the whole response distribution. A neuron can have the same mean response on two trials but different spike times, or the same mean tuning curve as another neuron but larger trial-to-trial variability. Neural coding begins by separating the conditional mean \(f_i(s)\) from the conditional distribution \(p(r_i\mid s)\).

\paragraph{Formal setup.}
Let \((\Omega,\mathcal F,P)\) be the probability space for repeated experimental trials. Let
\[
S:\Omega\to\mathcal S
\]
be a stimulus or task variable, and let
\[
R_i:\Omega\to\mathcal R_i
\]
be the response of neuron \(i\). For spike counts in a window, \(\mathcal R_i=\{0,1,2,\ldots\}\). For a smoothed firing rate, \(\mathcal R_i\subseteq\mathbb R_{\geq0}\).

The tuning curve is the conditional expectation
\[
\boxed{
f_i(s)=\mathbb E[R_i\mid S=s].
}
\]
The response variability around the tuning curve is described by
\[
\operatorname{Var}(R_i\mid S=s)
=\mathbb E[(R_i-f_i(s))^2\mid S=s],
\]
or more completely by the conditional law \(p(r_i\mid s)\).

A common count model is the Poisson tuning model
\[
R_i\mid S=s\sim \operatorname{Poisson}(f_i(s)\Delta t),
\]
where \(f_i(s)\) is a rate in spikes per second and \(\Delta t\) is the window length in seconds. The mean count is \(f_i(s)\Delta t\).

\paragraph{Core intuition.}
Think of a tuning curve as a conditional average drawn over the stimulus axis. The stimulus value chooses a slice of trials. Averaging responses inside that slice gives one point on the curve. Repeating this for many stimulus values gives the function \(f_i\).

The curve answers an encoding question: if the outside world or task variable were \(s\), what response should we expect from this neuron? It does not, by itself, answer the decoding question: after observing a response \(r_i\), what stimulus was present? Decoding also needs the prior over stimuli and the response variability.

\paragraph{Main mechanism: from response distributions to tuning curves.}
For a discrete response variable,
\[
f_i(s)=\sum_{r\in\mathcal R_i} r\,p(r\mid s).
\]
For a continuous response variable,
\[
f_i(s)=\int_{\mathcal R_i} r\,p(r\mid s)\,dr.
\]
Thus the tuning curve is a lossy summary of the conditional distribution. Two response models may share the same tuning curve but have different reliability. For example,
\[
R_i\mid s\sim\operatorname{Poisson}(f_i(s))
\]
has conditional variance \(f_i(s)\), while a Gaussian approximation
\[
R_i\mid s\sim\mathcal N(f_i(s),\sigma^2)
\]
has constant conditional variance \(\sigma^2\). The same conditional mean can support different uncertainty calculations.

For a differentiable scalar stimulus \(s\) and independent Poisson neurons with mean counts \(\lambda_i(s)\), the local Fisher information is
\[
\boxed{
\mathcal I(s)=\sum_{i=1}^N \frac{(\lambda_i'(s))^2}{\lambda_i(s)}.
}
\]
This follows from Chapter 18's Fisher-information definition applied to
\[
\log p(r\mid s)=\sum_i \left[r_i\log \lambda_i(s)-\lambda_i(s)-\log(r_i!)\right].
\]
Differentiating with respect to \(s\) and taking the response expectation leaves the squared sensitivity divided by the mean. Steep tuning curves help local discrimination, but only when response variability is not too large.

\paragraph{Worked examples.}
\emph{Small numerical example.}
An orientation-tuned neuron is tested at four orientations. Across repeated trials, the mean spike counts in a \(100\) ms window are
\[
\begin{array}{c|cccc}
s & 0^\circ & 45^\circ & 90^\circ & 135^\circ\\
\hline
f(s) & 1 & 4 & 8 & 4
\end{array}
\]
The preferred orientation among these tested values is \(90^\circ\). If the spike count at \(90^\circ\) has sample variance \(16\), the mean says the neuron is strongly driven, while the variance says individual trials are noisy.

\emph{Computational neuroscience example.}
In visual cortex, \(s\) may be grating orientation and \(R_i\) may be spike count after stimulus onset. In hippocampus, \(s\) may be animal location and \(f_i(s)\) is a place field. In motor cortex, \(s\) may be reach direction and \(f_i(s)\) may be direction tuned. In each case, the conditional mean is estimated by grouping trials or time bins by the variable \(s\).

\emph{AI example.}
A hidden unit activation in a neural network can be treated like a tuning curve over inputs. If \(S\) is an image attribute and \(R_i\) is activation of feature \(i\), then \(f_i(s)=\mathbb E[R_i\mid S=s]\) describes how that feature responds to the attribute. This does not prove that the artificial unit is biological, but it uses the same conditional-expectation mathematics.

\paragraph{Common misconceptions and failure modes.}
A tuning curve is not a causal proof that the neuron represents the stimulus. The response may covary with attention, movement, reward expectation, or another latent variable correlated with \(S\). A tuning curve is also not a single-trial decoder. High mean response to one stimulus does not imply that every high response came from that stimulus. Finally, a smooth plotted curve can hide finite-sample uncertainty; sparse stimulus sampling can make selectivity look sharper or weaker than it is.

\paragraph{Exercises.}
\begin{enumerate}[label=\textbf{47.1.\arabic*.}, leftmargin=*]
\item \textbf{Mechanical.} A neuron has mean spike counts \(2,5,3\) for stimuli \(s_1,s_2,s_3\). Write the tuning curve values and identify the preferred stimulus among the three.
\item \textbf{Proof or derivation.} Starting from \(f_i(s)=\mathbb E[R_i\mid S=s]\), derive the sum formula \(f_i(s)=\sum_r r\,p(r\mid s)\) for a discrete response.
\item \textbf{Numerical/computational.} Under a Poisson model with mean counts \(\lambda(0)=2\) and \(\lambda(1)=6\), compute \(P(R=4\mid S=0)\) and \(P(R=4\mid S=1)\). Which stimulus gives the larger likelihood for observing \(R=4\)?
\item \textbf{Interpretation.} Explain why a neuron with a clear tuning curve can still be unreliable on single trials.
\item \textbf{Debug the misconception.} A student says, "The tuning curve tells us exactly what stimulus occurred on each trial." Correct the statement using conditional means and response variability.
\end{enumerate}

\subsection{47.2 Population coding}

\paragraph{Motivating question.}
Why can a group of neurons encode a variable more precisely than any single neuron, and what geometry does the population response trace out?

\paragraph{Beginner setup.}
A population response is a vector. If \(N\) neurons are recorded in one time window, the response is
\[
R=(R_1,\ldots,R_N)\in\mathbb R^N
\]
or, for spike counts, \(R\in\mathbb N^N\). For each stimulus \(s\), the population has a mean response vector
\[
\mu(s)=\mathbb E[R\mid S=s]
=\bigl(f_1(s),\ldots,f_N(s)\bigr).
\]
As \(s\) changes, \(\mu(s)\) traces a curve or surface in the \(N\)-dimensional neural response space.

Population coding solves a distributed-representation problem. One neuron may be ambiguous because several stimuli produce similar responses. A population can disambiguate because different neurons have different tuning curves. The pattern across neurons matters.

\paragraph{Formal setup.}
Let \(S\in\mathcal S\) be a stimulus and let
\[
R=(R_1,\ldots,R_N):\Omega\to\mathcal R\subseteq\mathbb R^N
\]
be the population response. The population tuning map is
\[
\boxed{
\mu:\mathcal S\to\mathbb R^N,\qquad
\mu(s)=\mathbb E[R\mid S=s].
}
\]
The conditional noise covariance is
\[
\Sigma(s)=\operatorname{Cov}(R\mid S=s)
=\mathbb E[(R-\mu(s))(R-\mu(s))^\top\mid S=s].
\]
The off-diagonal entries of \(\Sigma(s)\) are noise correlations. They matter because two neurons with correlated trial-to-trial fluctuations may add less independent evidence than two neurons with independent fluctuations.

For a scalar stimulus \(s\) and approximately Gaussian population responses
\[
R\mid S=s\sim\mathcal N(\mu(s),\Sigma(s)),
\]
with \(\Sigma(s)=\Sigma\) independent of \(s\), the local Fisher information is
\[
\boxed{
\mathcal I(s)=\mu'(s)^\top \Sigma^{-1}\mu'(s).
}
\]
Here \(\mu'(s)\in\mathbb R^N\) is the tangent vector to the population response curve.

\paragraph{Core intuition.}
The population mean \(\mu(s)\) is a point in response space. A small change in stimulus moves that point in the direction \(\mu'(s)\). Noise covariance draws an uncertainty ellipsoid around it. Discrimination is easy when the stimulus-induced movement is large relative to the noise in that direction.

This is why population geometry matters. If \(\mu(s_1)\) and \(\mu(s_2)\) are far apart compared with noise, the population separates the stimuli. If they are close or if noise stretches along the same direction that separates them, decoding is harder.

\paragraph{Main mechanism: population-vector decoding.}
For circular variables such as direction, a simple population-vector decoder assigns each neuron a preferred unit vector \(v_i\in\mathbb R^2\). Given responses \(r_i\), define
\[
\boxed{
\widehat v=\sum_{i=1}^N r_i v_i,
\qquad
\widehat s=\operatorname{angle}(\widehat v).
}
\]
The decoder treats each neuron as voting for its preferred direction with weight proportional to response magnitude.

This method is exact only under special symmetric assumptions, but it teaches the central geometry: a scalar or circular variable can be represented by a distributed response vector, and decoding can be a projection back from response space to stimulus space.

\paragraph{Worked examples.}
\emph{Small numerical example.}
Four neurons prefer right, up, left, and down, with
\[
v_1=(1,0),\quad v_2=(0,1),\quad v_3=(-1,0),\quad v_4=(0,-1).
\]
If the observed responses are
\[
r=(8,4,1,3),
\]
then
\[
\widehat v=8(1,0)+4(0,1)+1(-1,0)+3(0,-1)=(7,1).
\]
The decoded direction is \(\operatorname{angle}(7,1)\), a small angle above right.

\emph{Computational neuroscience example.}
Motor-cortex population activity during reaching can be treated as a vector of firing rates. Even when individual neurons are broadly tuned, the population vector can predict reach direction. More modern analyses often replace this simple vector sum with probabilistic decoders, but the geometric idea remains: behavior is read from a pattern across neurons.

\emph{AI example.}
Word embeddings, image embeddings, and hidden-layer activations are population codes in artificial systems. A single coordinate rarely determines the represented object. Instead, semantic or task information is distributed across a vector. Distances and directions in the representation space become meaningful only after specifying the training distribution, loss, and readout.

\paragraph{Common misconceptions and failure modes.}
More neurons do not automatically mean more information. If many neurons are redundant or share strong noise correlations, adding them can provide little new evidence. A low-dimensional population curve is not necessarily the true stimulus space; it may reflect task design, preprocessing, or the chosen analysis window. Finally, a population vector decoder is a model, not a law of biology. It is useful when its assumptions approximate the data.

\paragraph{Exercises.}
\begin{enumerate}[label=\textbf{47.2.\arabic*.}, leftmargin=*]
\item \textbf{Mechanical.} For \(R=(2,0,5)\), write the response as a point in \(\mathbb R^3\). If \(\mu(s_1)=(1,1,1)\) and \(\mu(s_2)=(3,0,4)\), compute the Euclidean distance between the two mean responses.
\item \textbf{Proof or derivation.} For Gaussian responses with fixed covariance \(\Sigma\), derive the Fisher-information formula \(\mathcal I(s)=\mu'(s)^\top\Sigma^{-1}\mu'(s)\) from the score of \(\mathcal N(\mu(s),\Sigma)\).
\item \textbf{Numerical/computational.} Use the four-neuron population-vector rule to decode responses \(r=(1,7,2,1)\) for right, up, left, and down preferred directions.
\item \textbf{Interpretation.} Explain how positive noise correlation between two similarly tuned neurons can reduce the benefit of recording both.
\item \textbf{Debug the misconception.} A student says, "Population coding means each neuron encodes a different variable." Correct the statement using distributed representations.
\end{enumerate}

\subsection{47.3 Decoding and reconstruction}

\paragraph{Motivating question.}
Given a neural response, how can we estimate the stimulus, behavior, or latent variable that produced it?

\paragraph{Beginner setup.}
Encoding goes from variables to responses:
\[
p(r\mid s).
\]
Decoding goes from responses to estimates:
\[
g(r)=\widehat s
\quad\text{or}\quad
p(s\mid r).
\]
A decoder may output a class label, a continuous estimate, a reconstructed image, a reach direction, or a posterior distribution.

There are two major families. A Bayesian decoder starts from an encoding model and prior, then uses Bayes' theorem. A discriminative decoder learns a direct map from responses to targets, often by regression or classification. Both require evaluation on held-out data because a decoder can memorize noise in the training set.

\paragraph{Formal setup.}
Let \(S\in\mathcal S\) be a target variable and \(R\in\mathcal R\) a population response. A deterministic decoder is a function
\[
g:\mathcal R\to\widehat{\mathcal S}.
\]
For a loss function \(L(\widehat s,s)\), the risk of \(g\) is
\[
\mathcal R(g)=\mathbb E[L(g(R),S)].
\]
A Bayesian decoder computes
\[
\boxed{
p(s\mid r)=\frac{p(r\mid s)p(s)}{\sum_{s'}p(r\mid s')p(s')}
}
\]
for discrete stimuli, with an integral replacing the sum for continuous \(s\). The maximum a posteriori estimate is
\[
\widehat s_{\mathrm{MAP}}(r)=\operatorname*{arg\,max}_{s\in\mathcal S} p(s\mid r).
\]
For squared-error loss and a continuous stimulus, the Bayes action is the posterior mean
\[
\widehat s_{\mathrm{mean}}(r)=\mathbb E[S\mid R=r].
\]

A linear decoder for real-valued targets uses
\[
\widehat s=w^\top r+b,
\]
or for multiple outputs,
\[
\widehat y=Wr+b.
\]
The weights are usually fit by least squares, logistic regression, or another supervised objective.

\paragraph{Core intuition.}
Bayesian decoding asks which stimulus values make the observed response probable, while respecting the prior. A high response from a neuron tuned to \(90^\circ\) supports \(90^\circ\), but only in proportion to the response variability and prior probability of each orientation.

Linear decoding asks for a weighted combination of neural activities that predicts the target. It can work well even when the true encoding model is complicated, but it gives a readout model, not necessarily the exact computation performed by the brain.

\paragraph{Main mechanism: independent Poisson Bayesian decoder.}
Assume \(N\) conditionally independent spike counts given stimulus \(s\):
\[
R_i\mid S=s\sim\operatorname{Poisson}(\lambda_i(s)).
\]
Then
\[
p(r\mid s)=\prod_{i=1}^N e^{-\lambda_i(s)}\frac{\lambda_i(s)^{r_i}}{r_i!}.
\]
The log posterior is, up to terms not depending on \(s\),
\[
\log p(s\mid r)
=\log p(s)+\sum_{i=1}^N
\left[r_i\log\lambda_i(s)-\lambda_i(s)\right]+C(r).
\]
Therefore the MAP decoder chooses the stimulus that maximizes
\[
\boxed{
\log p(s)+\sum_{i=1}^N
\left[r_i\log\lambda_i(s)-\lambda_i(s)\right].
}
\]
The factorial terms \(r_i!\) cancel because the observed response \(r\) is fixed while comparing stimuli.

\paragraph{Worked examples.}
\emph{Small numerical example.}
Let \(S\in\{\mathrm{left},\mathrm{right}\}\) with equal prior. A population has one spike-count response \(R\). Suppose
\[
R\mid S=\mathrm{left}\sim\operatorname{Poisson}(2),
\qquad
R\mid S=\mathrm{right}\sim\operatorname{Poisson}(6).
\]
If \(R=5\), then
\[
P(R=5\mid S=\mathrm{right})=e^{-6}\frac{6^5}{5!},
\]
and
\[
P(R=5\mid S=\mathrm{left})=e^{-2}\frac{2^5}{5!}.
\]
Equal priors and common \(5!\) cancel, giving posterior probability
\[
P(S=\mathrm{right}\mid R=5)
=\frac{e^{-6}6^5}{e^{-2}2^5+e^{-6}6^5}
\approx 0.80.
\]

\emph{Computational neuroscience example.}
In retinal decoding, \(R\) may be a population spike pattern and \(S\) a visual stimulus frame or motion variable. A Bayesian decoder uses an encoding model \(p(r\mid s)\) to reconstruct likely stimuli. A linear decoder may instead learn filters that map binned spikes to stimulus pixels or motion components.

\emph{AI example.}
An encoder-decoder neural network maps an input to a latent representation and then decodes it into a reconstruction. A classifier layer is also a decoder: it maps representation \(r\) to label probabilities. The same evaluation principle applies as in neural decoding: performance must be measured on held-out data, not just training examples.

\paragraph{Common misconceptions and failure modes.}
A successful decoder does not prove that a downstream brain area uses the same algorithm. It proves that the information is available to that decoder class under the chosen data split. A poor decoder does not prove absence of information; the decoder may be misspecified or underpowered. Bayesian decoders are only as good as their encoding model and prior. Linear decoders are easy to interpret but may miss nonlinear information.

\paragraph{Exercises.}
\begin{enumerate}[label=\textbf{47.3.\arabic*.}, leftmargin=*]
\item \textbf{Mechanical.} Define the domain and codomain of a decoder that maps a \(100\)-neuron spike-count vector to one of eight reach directions.
\item \textbf{Proof or derivation.} Derive the independent Poisson MAP objective from the product likelihood and Bayes' theorem.
\item \textbf{Numerical/computational.} For priors \(P(s_1)=0.7\), \(P(s_2)=0.3\) and likelihoods \(p(r\mid s_1)=0.2\), \(p(r\mid s_2)=0.5\), compute the posterior over \(s_1,s_2\).
\item \textbf{Interpretation.} Why should decoder performance be evaluated on held-out trials or held-out time blocks?
\item \textbf{Debug the misconception.} A student says, "If a linear decoder fails, the neural population contains no information about the stimulus." Explain why the conclusion is too strong.
\end{enumerate}

\subsection{47.4 Information in neural responses}

\paragraph{Motivating question.}
How much uncertainty about a stimulus is reduced by observing a neural response, and how reliable are estimates of that quantity from finite data?

\paragraph{Beginner setup.}
Chapter 18 defined entropy, KL divergence, cross-entropy, mutual information, and data processing. Neural coding gives these objects concrete roles. The stimulus is a random variable \(S\). The response is a random variable \(R\), often a spike count, spike word, binned population vector, or latent response feature.

The total response entropy \(H(R)\) measures how variable responses are across all trials. The noise entropy \(H(R\mid S)\) measures response variability after the stimulus is fixed. The stimulus-response mutual information
\[
I(S;R)=H(R)-H(R\mid S)
\]
measures how much observing \(R\) reduces uncertainty about \(S\).

Information here means uncertainty reduction under a specified probability model. It does not mean semantic meaning, consciousness, importance, or causal influence by itself.

\paragraph{Formal setup.}
Assume first that \(S\) and \(R\) are discrete. Their joint distribution is
\[
p(s,r)=P(S=s,R=r),
\]
with marginals \(p(s)\), \(p(r)\), and conditional laws \(p(r\mid s)\), \(p(s\mid r)\).

The stimulus entropy is
\[
H(S)=-\sum_s p(s)\log_2 p(s).
\]
The response entropy is
\[
H(R)=-\sum_r p(r)\log_2 p(r).
\]
The conditional response entropy, often called noise entropy in coding experiments, is
\[
H(R\mid S)=-\sum_s p(s)\sum_r p(r\mid s)\log_2 p(r\mid s).
\]
The stimulus-response mutual information is
\[
\boxed{
I(S;R)=\sum_{s,r}p(s,r)\log_2\frac{p(s,r)}{p(s)p(r)}
=H(S)-H(S\mid R)=H(R)-H(R\mid S).
}
\]

If a fitted response model \(q(r\mid s)\) approximates the true \(p(r\mid s)\), the conditional cross-entropy is
\[
H(p,q\mid S)
=-\sum_s p(s)\sum_r p(r\mid s)\log_2 q(r\mid s),
\]
and the conditional KL mismatch is
\[
\sum_s p(s)D_{\mathrm{KL}}(p(r\mid s)\|q(r\mid s)).
\]
This connects neural encoding models to prediction loss.

\paragraph{Core intuition.}
Response entropy alone is not information about the stimulus. A neuron with high spontaneous variability can have high \(H(R)\) and low \(I(S;R)\). Information appears when stimulus identity changes the response distribution enough that observing the response reduces uncertainty about the stimulus.

The data-processing inequality gives a useful discipline. If
\[
S\to R\to Z
\]
is a Markov chain, where \(Z\) is a processed version of the response that receives no additional information about \(S\), then
\[
I(S;Z)\leq I(S;R).
\]
A decoder output, low-dimensional projection, or learned representation cannot contain more information about \(S\) than the response it was computed from, unless it uses additional inputs.

\paragraph{Main mechanism: total entropy minus noise entropy.}
Starting from mutual information,
\[
I(S;R)=\sum_{s,r}p(s,r)\log_2\frac{p(r\mid s)}{p(r)}.
\]
Expand:
\[
I(S;R)=
\sum_{s,r}p(s,r)\log_2 p(r\mid s)
-
\sum_{s,r}p(s,r)\log_2 p(r).
\]
The first term is \(-H(R\mid S)\). In the second term, \(\sum_s p(s,r)=p(r)\), so
\[
-\sum_{s,r}p(s,r)\log_2 p(r)
=-\sum_r p(r)\log_2 p(r)
=H(R).
\]
Thus
\[
\boxed{I(S;R)=H(R)-H(R\mid S).}
\]
The information is the part of response variability explained by stimulus variation, measured in bits.

\paragraph{Worked examples.}
\emph{Small joint-table example.}
Let \(S\in\{\mathrm{low},\mathrm{high}\}\) and \(R\in\{\mathrm{quiet},\mathrm{active}\}\), with
\[
\begin{array}{c|cc}
 & \mathrm{quiet} & \mathrm{active}\\
\hline
\mathrm{low} & 0.45 & 0.05\\
\mathrm{high} & 0.05 & 0.45
\end{array}
\]
Both marginals are uniform, so \(H(S)=H(R)=1\) bit. Given \(S=\mathrm{low}\), the response distribution is \((0.9,0.1)\), whose entropy is
\[
-0.9\log_2(0.9)-0.1\log_2(0.1)\approx 0.469.
\]
The same conditional entropy holds for \(S=\mathrm{high}\), so
\[
H(R\mid S)\approx 0.469,\qquad
I(S;R)\approx 1-0.469=0.531\ \text{bits}.
\]

\emph{Finite-data example.}
If \(R\) is binned into many response words but only a few trials are observed, the empirical distribution often assigns too many zero and one-count bins. The plug-in estimate of entropy and mutual information can be biased. Coarser bins reduce variance but may lose structure. Model-based estimates reduce dimensionality but introduce assumptions through \(q(r\mid s)\).

\emph{Computational neuroscience example.}
A spike train can be converted into a word by dividing a \(100\) ms window into ten \(10\) ms bins and marking each bin as spike or no spike. Then \(R\) has up to \(2^{10}\) possible words. Mutual information \(I(S;R)\) asks how much those words reduce stimulus uncertainty. The answer depends on bin width, trial count, response window, and stationarity.

\emph{AI example.}
Let \(X\) be an input image, \(Z=f_\theta(X)\) a representation, and \(Y\) a class label. Data-processing intuition says that deterministic processing cannot create new information about \(X\), but it can discard nuisance information while preserving label-relevant information. In practice, information estimates for high-dimensional neural-network activations require approximations or bounds, so numerical values should be interpreted carefully.

\paragraph{Common misconceptions and failure modes.}
Mutual information is not vague meaning. It is a functional of a joint distribution. It does not imply causality, and it does not identify the decoding algorithm used by the brain. Finite-data estimates can be severely biased, especially for high-dimensional responses, continuous variables, adaptive binning, or correlated trials. Comparing information across experiments requires matching stimulus ensembles and response windows.

\paragraph{Exercises.}
\begin{enumerate}[label=\textbf{47.4.\arabic*.}, leftmargin=*]
\item \textbf{Mechanical.} If \(H(R)=3.0\) bits and \(H(R\mid S)=1.2\) bits, compute \(I(S;R)\). Interpret the result as uncertainty reduction.
\item \textbf{Proof or derivation.} Starting from \(I(S;R)=D_{\mathrm{KL}}(p(s,r)\|p(s)p(r))\), derive \(I(S;R)=H(S)-H(S\mid R)\).
\item \textbf{Numerical/computational.} For the joint table in the worked example, change the diagonal entries to \(0.4\) and off-diagonal entries to \(0.1\). Compute \(I(S;R)\).
\item \textbf{Interpretation.} A response has high entropy but zero mutual information with the stimulus. What does that imply about \(p(r\mid s)\) across stimuli?
\item \textbf{Debug the misconception.} A student says, "Information is whatever the response means to the animal." Correct the statement using the joint distribution \(p(s,r)\) and uncertainty reduction.
\end{enumerate}

\paragraph{Closing bridge.}
Neural coding begins with conditional response distributions, tuning curves, and decoders, then uses information theory to quantify uncertainty reduction. Chapter 48 keeps the population-response viewpoint but shifts from coding variables to geometry and dynamics: activity vectors become points, trajectories, manifolds, and latent states inferred from high-dimensional neural data.

\clearpage
